% Generated draft: requires mathematical and source review.
\documentclass[11pt]{article}
\usepackage[margin=25mm]{geometry}
\usepackage{fontspec}
\setmainfont{Latin Modern Roman}
\usepackage{amsmath,amssymb,mathtools,unicode-math}
\setmathfont{Latin Modern Math}
\usepackage{enumitem,needspace,xurl,hyperref}
\setlength{\parindent}{0pt}
\setlength{\parskip}{5pt}
\setlength{\emergencystretch}{3em}
\newcommand{\N}{\mathbb{N}}
\newcommand{\Z}{\mathbb{Z}}
\newcommand{\Q}{\mathbb{Q}}
\newcommand{\R}{\mathbb{R}}
\newcommand{\C}{\mathbb{C}}
\newcommand{\F}{\mathbb{F}}
\newcommand{\A}{\mathbb{A}}
\newcommand{\PP}{\mathbb{P}}
\newcommand{\E}{\mathbb{E}}
\newcommand{\eps}{\varepsilon}
\newcommand{\dd}{\,\mathrm{d}}
\begin{document}
\section*{UnsolvedMath ID: 20003352}
\textbf{Source problem code:} catalogue entry
\subsection{Definitions and mathematical statement}
This is an expository motivation paragraph following Conjecture 51, not a separate conjecture. The source's truncated final phrase is completed as ``a proof of Conjecture 51'' according to the supplied aid. The antecedent, recovered from companion permanent ID $20003350$, is the Andr\'e--Oort assertion: the Zariski closure of a set of special points in a Shimura variety is a finite union of special subvarieties.

The paragraph draws an analogy with an abelian variety $A$: special points correspond to torsion points $P$ with $nP=0$ for some $n\ge1$, and special subvarieties correspond to torsion translates $P+B$ of abelian subvarieties $B\subset A$. It attributes the analogous conjectures to S. Lang, Yu. Manin and D. Mumford, reports them settled with survey [190], and describes proofs [191, 208] using equidistribution of Galois orbits of suitable ``generic'' sequences [186], also [20]. It does not state the precise height and genericity hypotheses of those equidistribution results; it is not an assertion about unrestricted sequences of algebraic points.

The attached aid reports this fixed-field conditional reduction. Let $Z$ be an irreducible algebraic subvariety of a Shimura variety with Zariski-dense special points, and let $M$ be its irreducible minimal special envelope: $Z\subset M$, with $Z$ not contained in any proper special subvariety of $M$. Assume one number field $E$ defines both $Z$ and $M$ and their ambient model. Suppose the proper special subvarieties of $M$ form a countable family and that every strict sequence of special points $p_n\in M(\overline E)$ has its $E$-Galois orbit probabilities
\[
 \mu_{p_n,E}=\frac1{|\operatorname{Gal}(\overline E/E)\cdot p_n|}
 \sum_{p\in\operatorname{Gal}(\overline E/E)\cdot p_n}\delta_p
\]
converging weakly as probabilities to a canonical measure $\mu_M$ with full support on $M(\C)$. Here strict means that every proper special subvariety of $M$ contains only finitely many $p_n$. Then the reported conclusion is
\[
 Z=M,
\]
so $Z$ is special.

The role of the hypotheses is explicit: countable diagonal avoidance produces a strict sequence from the dense special points of $Z$; the single field of definition $E$ makes every orbit measure supported on the same closed set $Z(\C)$. Portmanteau then gives $\mu_M(Z(\C))=1$, and full support forces $Z=M$. No changing sequence of defining fields is permitted. The full-support orbit-equidistribution hypothesis is the substantive conditional input, not proved here or inferred from the analogy alone. The assigned registry record labels this material context only; its broader neighboring status claims do not turn this reduction into an unconditional new proof.
\subsection{Short English statement}
The paragraph motivates Andr\'e--Oort by the torsion-point analogy for abelian varieties. Its supplied reduction is conditional: equidistribution of strict special-point Galois orbits over one fixed defining field, with full support on the minimal special envelope, forces a subvariety with dense special points to be special.
\subsection{Sources}
\begin{itemize}
\item American Institute of Mathematics, \emph{Emerging applications of measure rigidity}, motivation following Conjecture 51, \url{https://aimath.org/WWN/measrigid/measrigid.pdf}. Source attributions: S. Lang, Yu. Manin and D. Mumford; survey [190]; proofs [191, 208]; equidistribution [186], also [20]. Complete bibliography entries are absent.
\item Companion local record, permanent ID $20003350$, source code \texttt{AIM-OTHER-0157}, canonical \texttt{aim-other-notes.json}, \texttt{notes[156]}, states Conjecture 51 attributed to Y. Andr\'e and F. Oort [3, 153]. Its aid distinguishes published conditional results, unconditional $\mathcal A_g$ results, and the general Pila--Shankar--Tsimerman preprint \texttt{arXiv:2109.08788v4}; these are supplied status reports, not independently checked here.
\item Local UnsolvedMath registry, permanent ID $20003352$, source code \texttt{catalogue entry}, canonical \texttt{aim-other-notes.json}, \texttt{notes[158]}, dated 14 August 2026. Research attempt 1: context only, ``not a problem'' at run, proved rigor, low novelty confidence; machine-generated and not peer reviewed.
\end{itemize}
\end{document}
